\documentclass[12pt,a4paper]{article}
\usepackage[UTF8]{ctex}
\usepackage{amsmath,amssymb,bm}
\usepackage{geometry}
\usepackage{booktabs}
\geometry{margin=2.5cm}

\title{四节点四边形等参单元分片试验}
\author{}
\date{}

\begin{document}

\maketitle

\section{题目}

利用 STAPpy/elasticity2d-python/ABAQUS 对四节点四边形等参单元进行分片试验。考虑平面应力问题，取
\begin{equation}
E = 1000, \qquad \nu = 0.3.
\end{equation}

给定位移场为
\begin{equation}
u(x,y) = 0.01x, \qquad v(x,y) = -0.003y.
\end{equation}

验证该位移场所对应的应力状态，并说明四节点四边形等参单元是否通过分片试验。

\section{理论分析}

在线弹性小变形条件下，平面应力问题的应变分量为
\begin{equation}
\varepsilon_x = \frac{\partial u}{\partial x},
\qquad
\varepsilon_y = \frac{\partial v}{\partial y},
\qquad
\gamma_{xy} = \frac{\partial u}{\partial y} + \frac{\partial v}{\partial x}.
\end{equation}

代入题设位移场可得
\begin{equation}
\varepsilon_x = 0.01,
\qquad
\varepsilon_y = -0.003,
\qquad
\gamma_{xy} = 0.
\end{equation}

故应变向量为
\begin{equation}
\bm{\varepsilon}=
\begin{bmatrix}
0.01 \\
-0.003 \\
0
\end{bmatrix}.
\end{equation}

平面应力本构矩阵为
\begin{equation}
\bm{D} = \frac{E}{1-\nu^2}
\begin{bmatrix}
1 & \nu & 0 \\
\nu & 1 & 0 \\
0 & 0 & \dfrac{1-\nu}{2}
\end{bmatrix}.
\end{equation}

将 $E=1000$、$\nu=0.3$ 代入，有
\begin{equation}
\bm{\sigma} = \bm{D}\bm{\varepsilon}.
\end{equation}

逐项计算得
\begin{align}
\sigma_x &= \frac{1000}{1-0.3^2}(0.01+0.3\times(-0.003)) = 10, \\
\sigma_y &= \frac{1000}{1-0.3^2}(0.3\times0.01-0.003) = 0, \\
\tau_{xy} &= \frac{1000}{2(1+0.3)}\times 0 = 0.
\end{align}

因此该位移场对应的精确应力解为
\begin{equation}
\sigma_x = 10,
\qquad
\sigma_y = 0,
\qquad
\tau_{xy} = 0.
\end{equation}

同时，体力项为
\begin{equation}
\bm{b} = \bm{0}.
\end{equation}

这说明该问题对应一个常应变、常应力的单向拉伸状态，是分片试验的标准检验情形。

\section{数值验证方法}

由于当前 STAP++ 仓库仅实现了杆单元，并未实现 Q4 单元，因此单独编写了 Python 脚本 STAPpy test.py，对四节点四边形双线性等参单元进行验证。数值模型采用畸变的 $2\times2$ 网格，共 4 个 Q4 单元、9 个节点，在边界节点施加精确位移，只保留中心节点为自由节点，采用 $2\times2$ Gauss 积分计算单元刚度矩阵与高斯点应力。

若单元通过分片试验，则应满足以下两点：
\begin{enumerate}
  \item 自由节点的数值位移与精确位移一致；
  \item 各高斯积分点的应力与理论常应力解一致。
\end{enumerate}

\section{数值结果}

Python 脚本输出结果为
\begin{verbatim}
Q4 patch test on a distorted 2x2 mesh
E = 1000.0, nu = 0.3
Exact strain = [ 0.01  -0.003  0.   ]
Expected stress = [1.0000000e+01 1.5335718e-17 0.0000000e+00]
Solved center displacement = [ 0.01  -0.003]
Exact center displacement = [ 0.01  -0.003]
Max center displacement error = 1.734723e-18
Max Gauss-point stress error = 5.329071e-15
Max free-DOF residual = 8.881784e-16
Patch test passed.
\end{verbatim}

由结果可见：
\begin{equation}
\max |\Delta \bm{u}| \approx 1.73\times 10^{-18},
\qquad
\max |\Delta \bm{\sigma}| \approx 5.33\times 10^{-15}.
\end{equation}

误差仅为机器舍入误差量级，可认为数值解与理论解完全一致。

\section{结论}

对于题设位移场
\begin{equation}
u(x,y) = 0.01x, \qquad v(x,y) = -0.003y,
\end{equation}
四节点四边形双线性等参单元能够精确再现线性位移场，并得到常应力解
\begin{equation}
\sigma_x = 10,
\qquad
\sigma_y = 0,
\qquad
\tau_{xy} = 0.
\end{equation}

因此，四节点四边形等参单元在该平面应力问题下通过了分片试验。

\end{document}